\subsection{实验结果分析}
\begin{itemize}
    \item 单线程算法性能分析:
    \begin{itemize}
        \item 现象:随着矩阵规模$N$的增大，所有算法的执行时间呈指数级增长。在小规模（32-128）下，各算法差距较小；但在大规模（512-1024）下，Naive (i-j-k) 版本的斜率明显高于其他版本，而 IKJ 和 Loop Unrolling 版本表现出明显的性能优势
        \item 原因:Naive 版本在最内层循环访问$B$矩阵时跨度大,空间局部性差,导致严重的Cache Miss。IKJ通过交换循环次序,使最内层循环对矩阵$B$的访问变为连续内存空间，极大提高了空间局部性，程序运行效率大大提升。
    \end{itemize}
    \item 多线并行与加速比分析:
    \begin{itemize}
        \item 现象:
        \begin{itemize}
            \item 大计算规模优势:当规模达到$1024 \times 1024$时,随着线程数增加,加速比呈现明显的上升趋势。
            \item 小规模"反向优化":在$N = 32$或$64$时,增加线程数不仅没有提速,加速比反而小于1(即变慢)。
            \item QEMU瓶颈:在线程数接近或超过8时(-smp 8),加速比曲线趋于平缓,甚至在$T = 12$时略有下滑。
            \item WSL线程与核心数权衡:WSl中当线程数为12(大于核心数6)时,运行效率低于线程数为6的情况。
        \end{itemize}
        \item 原因:
        \begin{itemize}
            \item 线程开销控制:小规模任务中,计算任务本身极快(微秒级),而多线程的创建、同步和销毁开销远超计算量,导致效率下降，此时多线程体现不出优势。
            \item 模拟器资源争抢：QEMU或WSL2 Arch linux 在宿主机上运行，当模拟线程数超过宿主机物理核心（6核）或分配核心（8核）时，会产生剧烈的上下文切换开销，导致系统负载过载。
        \end{itemize}
    \end{itemize}
    \item 仿真环境与真实环境差异:
    \begin{itemize}
        \item 现象:WSL宿主机上的运行时间比QEMU低约2各数量级;WSL 环境下的曲线更为平稳且符合理论预期;QEMU 下的曲线存在细微波动，且在处理超大规模矩阵时会出现系统调用报错或计时器延迟。
        \item 原因:QEMU 作为系统级模拟器，需要将每一条 RISC-V 指令动态翻译为 x86 指令执行，这带来了巨大的软件模拟损耗。同时宿主机拥有真实的硬件缓存和预取器（L1d 288KB），而 QEMU 的缓存层级是软件定义的，无法完美重现硬件级别的访存延迟。
    \end{itemize}
    \item 其他情况分析:
    \begin{itemize}
        \item 现象
        \begin{itemize}
            \item 矩阵转置优化效果并不是很好。
            \item 矩阵分块算法优化随着分块的增大,程序性能可能会降低。
        \end{itemize}
        \item 原因
        \begin{itemize}
            \item 矩阵转置优化版的代码要先执行矩阵优化,这一步会拖慢矩阵乘法的计算,即使之后能使得矩阵$B$的访问也是内存连续访问。
            \item 当矩阵分块大小超过L1/L2cache时,会使得即使是分块计算矩阵乘法也会出现较大的Cache Miss率,且矩阵分块乘法算法本身较为复杂,所以实际使用时分块优化的效果也不是最好的。
        \end{itemize}
        
    \end{itemize}

\end{itemize}
